\documentclass[11pt]{article}
\usepackage[margin=1in]{geometry}
\usepackage{amsmath}
\usepackage{amssymb}
\usepackage{enumitem}
\usepackage{titlesec}
\usepackage{microtype}

\titleformat{\section}{\Large\bfseries}{\thesection}{1em}{}[\titlerule]
\titleformat{\subsection}{\large\bfseries}{\thesubsection}{1em}{}

\newcommand{\coursestatus}[1]{\hfill\normalsize\normalfont\textit{#1}}

\title{\textbf{Graduate Mathematics Coursework and Academic Background}}
\author{Bikash Thakur}
\date{}

\begin{document}

\maketitle

\section{Analysis}

\subsection{MTH 541. Real Analysis \mdseries(3 credits)}
Continuity, differentiation, convergence, series and integration, in both one and several variables.

\subsection{MTH 641. Functions of a Real Variable \mdseries(4 credits)}
Lebesgue measure, Lebesgue integration, differentiation, general measures and integration, Radon--Nikodym theorem, Fubini theorem, classical $L^p$ spaces, Banach spaces.

\subsection{MATH 645. Analysis I \mdseries(3 credits)}
Review and extension of the fundamental concepts of advanced calculus: the real number system, limit, continuity, differentiation, the Riemann integral, sequences and series. Point set topology in metric spaces. Uniform convergence and its applications.

\subsection{MATH 745. Analysis II \mdseries(3 credits)}
Lebesgue measure and integration, including the Lebesgue dominated convergence theorem and Riesz--Fischer theorem. Elements of Hilbert spaces and $L^p$-spaces. Fourier series and harmonic analysis. Multivariate calculus.

\subsection{MATH 656. Complex Variables I \mdseries(3 credits)}
The theory and applications of analytic functions of one complex variable: elementary properties of complex numbers, analytic functions, elementary complex functions, conformal mapping, Cauchy integral formula, maximum modulus principle, Laurent series, classification of isolated singularities, residue theorem, and applications.

\section{Algebra}

\subsection{MTH 521. Introduction to Abstract Algebra \mdseries(4 credits)}
Elementary theory of groups, rings, integral domains, fields, homomorphisms, and quotient structures.

\subsection{MTH 621. Abstract Algebra I \mdseries(4 credits)}
Sylow theory, composition series, polynomial rings. Galois theory of fields, modules over a principal ideal domain and their application.

\section{Linear Algebra and Matrix Theory}

\subsection{MTH 533. Applied Linear Algebra \mdseries(3 credits)}
A course in linear algebra with a focus on applications and implementation of those applications using current computational software. Topics such as singular value decomposition, matrix factorizations, stochastic matrices and eigenvalue approximation applied to problems in spline fitting, principal component analysis, random walks, image processing, least squares and recommender systems.

\subsection{MATH 631. Linear Algebra \mdseries(3 credits)}
Linear systems of equations, matrix algebra, linear spaces, orthogonality, eigenvalues and eigenvectors, diagonalization, and matrix decomposition. Applications.

\subsection{MATH 707. Numerical Linear Algebra \mdseries(3 credits)}
Algorithms and analysis for the fundamental numerical methods of linear algebra. QR factorizations, least squares problems, conditioning, stability, eigenvalue computations, and iterative methods.

\section{Topology}

\subsection{MTH 591. Introduction to Topology \mdseries(3 credits)}
Topological spaces, continuity, countability and separation axioms, product topology, quotient topology, compactness, connectedness, Tychonoff's Theorem, Urysohn's Lemma, metrizable spaces.

\subsection{MTH 691. Introduction to Algebraic Topology \mdseries(4 credits)}
Metric spaces, continuity, topological spaces, sums, products and quotients of topological spaces, compact spaces, connected spaces. Additional topics include either advanced topics in set-theoretic topology (including separation axioms, locally compact spaces, convergence, compactness and countability properties, Stone--\v{C}ech compactification, paracompact spaces, and metrization) or advanced topics in algebraic topology (including the fundamental group, The Seifert--Van Kampen Theorem, group presentations, covering spaces, the correspondence theorem, and topics in homology).

\section{Differential Equations}

\subsection{MTH 655. Advanced Differential Equations \mdseries(3 credits)}
Concepts and techniques for solving the ordinary and partial differential equations that arise in various scientific disciplines.

\subsection{MATH 651. Methods of Applied Mathematics I \mdseries(3 credits)}
A survey of mathematical methods for the solution of problems in the applied sciences and engineering. Topics include: ordinary differential equations and elementary partial differential equations. Fourier series, Fourier and Laplace transforms, and eigenfunction expansions.

\subsection{MATH 689. Advanced Applied Mathematics II: Ordinary Differential Equations \mdseries(3 credits)}
A practical and theoretical treatment of boundary-value problems for ordinary differential equations: generalized functions, Green's functions, spectral theory, variational principles, and allied numerical procedures. Examples drawn from applications in science and engineering.

\subsection{MATH 676. Advanced Ordinary Differential Equations \mdseries(3 credits) \coursestatus{Fall 2026}} 
A rigorous treatment of the theory of systems of differential equations: existence and uniqueness of solutions, dependence on initial conditions and parameters. Linear systems, stability, and asymptotic behavior of solutions. Nonlinear systems, perturbation of periodic solutions, and geometric theory of systems of ODEs.

\subsection{MATH 690. Advanced Applied Mathematics III: Partial Differential Equations \mdseries(3 credits) \coursestatus{Fall 2026}}
A practical and theoretical treatment of initial- and boundary-value problems for partial differential equations: Green's functions, spectral theory, variational principles, transform methods, and allied numerical procedures. Examples drawn from applications in science and engineering.

\section{Optimization}

\subsection{MTH 532. Optimization \mdseries(3 credits)}
Optimization of functions of several variables, convexity and least squares, Kuhn--Tucker conditions, linear programming.

\subsection{MTH 632. Advanced Optimization \mdseries(3 credits)}
Careful development of the theory of finite-dimensional continuous optimization, emphasizing the differentiable and convex cases.

\subsection{MATH 707. Optimization \mdseries(3 credits) \coursestatus{Audit Fall 2025}}
Modern topics in continuous optimization. Convex analysis including conditions for optimality, systematic development of interior-point methods, and modern applications.

\section{Graph Theory}

\subsection{MTH 538. Theory and Applications of Graphs \mdseries(3 credits) \coursestatus{Audit Fall 2024}}
Basic structural properties of graphs, trees, connectivity, traversability (Eulerian Tours and Hamiltonian Cycles), matchings, and vertex and edge colorings. Classic graph algorithms including shortest path, minimum weight tree, optimal assignment. Additional topics from network flows, planarity, extremal problems, and directed graphs.

\subsection{MTH 638. Advanced Graph Theory \mdseries(3 credits) \coursestatus{Audit Spring 2025}}
Advanced treatment of graph theory with selected topics from: Extremal problems, probabilistic, algebraic, and topological aspects of graph theory, analysis of graph algorithms, Ramsey theory.

\section{Numerical Analysis and Computational Mathematics}

\subsection{MTH 553. Numerical Analysis \mdseries(3 credits)}
Errors and error propagation, root-finding methods, numerical solution of linear systems, polynomial and cubic spline interpolation, numerical differentiation and integration, programming of algorithms.

\subsection{MATH 614. Numerical Methods I \mdseries(3 credits)}
Theory and techniques of scientific computation. Machine arithmetic, numerical solution of linear systems and pivoting, interpolation and quadrature, iterative solution of nonlinear systems, computation of eigenvalues and eigenvectors, numerical solution of initial- and boundary-value problems for systems of ODEs. Applications.

\subsection{MATH 712. Numerical Methods II \mdseries(3 credits) \coursestatus{Fall 2026}}
Numerical methods for the solution of initial- and boundary-value problems for partial differential equations, with emphasis on finite difference methods. Consistency, stability, convergence, and implementation.

\section{Applied Mathematics and Modeling}

\subsection{MTH 547. Topics in Mathematical Finance \mdseries(3 credits)}
Mathematical methods in options pricing; options and their combinations, arbitrage and put-call parity, stock and option trees, risk neutral pricing, geometric Brownian motion for stock models and derivation of the Black--Scholes formula. Additional topics: futures, forwards, swaps and bond models.

\subsection{MATH 613. Advanced Applied Mathematics I: Modeling \mdseries(3 credits)}
Concepts and strategies of mathematical modeling. Consistency of a model, nondimensionalization and scaling, regular and singular effects. Possible topics include continuum mechanics (heat and mass transfer, fluid dynamics, elasticity), vibrating strings, population dynamics, traffic flow, and the Sommerfeld problem.

\section{Physics}

\subsection{PHY 611. Advanced Classical Mechanics \mdseries(3 credits) \coursestatus{Fall 2026}}
General formalism of Lagrangian and Hamiltonian mechanics. D'Alembert's principle and virtual work problems. Euler--Lagrange equation for mechanics problems, constraint forces, and Lagrangian Multiplier method. Central force problems and the Kepler problem. Rigid body rotation using Euler's equation. Small oscillation and eigenvalue problems. Variational principle and its equivalence to the Lagrangian Equation. Canonical transformations using Poisson bracket methods. Hamilton--Jacobi equation and action-angle variables for periodic problems.

\section{Teaching Experience}

\subsection{MTH 151. Calculus I \mdseries(4 credits) --- Instructor of Record}
Topics include limits and continuity, derivatives and their applications, and early integration techniques of polynomial, rational, radical, trigonometric, inverse trigonometric, exponential, and logarithmic functions.

\section{Financial Engineering (WorldQuant University)}

\subsection{MScFE 560. Financial Markets \mdseries(4 credits)}
Introduction to professional finance: markets, products, participants, and regulation. Activities within financial markets including trading, financing, brokering, pricing, hedging, optimizing, and managing risk. Exploration of asset classes including rates, equities, cryptocurrencies, credit products, derivatives, ETFs, securitized products, real estate, commodities, and FX.

\subsection{MScFE 600. Financial Data \mdseries(4 credits)}
Empirical data for decision making and predictions in finance. Data sets including traditional time series, macroeconomics, and alternative data (text, social media, images, geolocation, climate). Application of Python for data selection, import, filtering, structuring, visualization, and analysis. Preparation of numeric and non-numeric data for financial models. Matrix transformations for normalization, reduction, factorization, and interpolation.

\subsection{MScFE 610. Financial Econometrics \mdseries(4 credits)}
Modeling probability distributions of returns using graphical, Bayesian, and non-parametrical methods. Univariate time series modeling: moving averages, autocorrelations, volatilities, GARCH models. Correlation, vector autoregressions, and cointegration. Statistical foundation and Python coding for econometric models. Application of bias, variance, and overfitting concepts to machine learning.

\subsection{MScFE 620. Derivative Pricing \mdseries(4 credits)}
Pricing options using stochastic calculus. No-arbitrage and perfect replication concepts. Black--Scholes Model. Construction of pricing models: binomial trees and finite difference methods for vanilla and exotic options. Price sensitivities (Greeks). Extensions including Heston Model and jump models. Python implementations.

\subsection{MScFE 622. Stochastic Modeling \mdseries(4 credits)}
Advanced volatility models upgrading Black--Scholes parameters to variables. Jump diffusions and heteroskedasticity. Markov processes including hidden Markov processes and Markov decision processes. Mathematical foundations for deep learning. Evaluation of assumptions, benefits, and difficulties of stochastic models.

\subsection{MScFE 632. Machine Learning in Finance \mdseries(4 credits)}
Fundamentals of supervised and unsupervised machine learning. Mathematical and computational foundations using Python and TensorFlow. Application to prediction, explanation, and comparison of financial series. Detection of overfitting using regularization, cross-validation, and resampling. Theoretical and practical background in deep learning models.

\subsection{MScFE 642. Deep Learning for Finance \mdseries(4 credits)}
Neural networks: CNNs (Convolutional Neural Networks), RNNs (Recurrent Neural Networks), LSTMs (Long Short-Term Memory), GRUs (Gated Recurring Units). Hyperparameter tuning using classical, Bayesian, and stochastic methods. Applications: stock price prediction, investment factor discovery, trading strategy back-testing. Transfer learning and data augmentation. Comprehensive intraday trading strategy development.

\subsection{MScFE 652. Portfolio Management \mdseries(4 credits) \coursestatus{Yet to Take}}
Portfolio optimization methodologies. Classical portfolio theory including Markowitz optimization. Modern approaches: Black--Litterman, probabilistic scenario optimization, prospect theory, Kelly criterion, risk parity. Integration of advanced econometrics and machine learning: neural networks, genetic algorithms, information theory, reinforcement learning. Applications across multiple asset classes.

\subsection{MScFE 660. Risk Management \mdseries(4 credits) \coursestatus{Yet to Take}}
Classical and modern methods of modeling and managing risk. Metrics and models for market, credit, and systemic risk. Integration of machine learning with classical methods (VaR, GARCH) and robust methods (Extreme Value Theory). Bayesian methods and Bayesian networks for systemic risk modeling. Complex network and scenario analysis for portfolio and systemic risk.

\subsection{MScFE 690. Capstone Project \mdseries(4 credits) \coursestatus{Yet to Complete}}
Application of financial engineering knowledge through project milestones. Development of problem statement, identification of required technology, peer review iterations, and final presentation. Options: Research track (investigating market inefficiency, risk management, alpha generation, arbitrage, asset pricing) or Practical track (development of operational trading systems, alpha generation modules, execution strategies). Deliverables: article-format paper and source code.

\section{Undergraduate Education (St. John's College, Santa Fe)}

\textit{Four-year all-required Great Books curriculum with no electives or majors. The integrated program yields substantial majors in History of Mathematics and Science and Philosophy, with minors in Classical Studies and Comparative Literature. Total: 136--137 semester hours.}

\subsection{Mathematics Curriculum \mdseries(36 credits)}
\begin{itemize}[noitemsep, topsep=0pt]
    \item \textbf{Geometry and Astronomy:} Euclid's Elements (synthetic plane and solid geometry, logical foundations, axiomatic theory of ratios); Ptolemy's Almagest (plane and spherical trigonometry, astronomical observations, mathematical treatment of planetary motion); Apollonius' Conics (geometrical exposition of conic curves).
    \item \textbf{Analytic Geometry and Calculus:} Analytic geometry of straight lines, circles, and conic sections (Descartes, Pascal); Calculus: limit theory, derivatives, integrals, logarithmic functions, differential equations (Galileo, Leibniz, Newton).
    \item \textbf{Physics and Foundations:} Newton's Principia Mathematica (mathematical treatment of space, time, force, and planetary celestial mechanics); Foundations of number theory (Dedekind); Non-Euclidean geometry (Lobachevsky); Special theory of relativity (Einstein).
\end{itemize}

\subsection{Laboratory Science Curriculum \mdseries(28.5 credits)}
\begin{itemize}[noitemsep, topsep=0pt]
    \item \textbf{Biology:} Observational biology, taxonomy, morphology, plant and animal functions (Aristotle, Bernard, Darwin, Driesch, Galen, Harvey, Linnaeus, Mendel, Morgan, Watson and Crick). Evolutionary theory, genetics, molecular biology.
    \item \textbf{Chemistry:} Concepts of weight, pressure, heat, temperature, gas laws (Archimedes, Pascal). Atomic theory (Avogadro, Cannizzaro, Dalton, Gay-Lussac, Lavoisier, Mendeleyev, Priestley).
    \item \textbf{Physics:} Classical kinematics and dynamics with attention to periodic motions (Descartes, Galileo, Leibniz, Newton). Thermodynamics. Electricity and magnetism culminating in Maxwell's equations (Faraday, Maxwell). Optics and mechanics. Atomic theory, statistical mechanics, quantum theory (Bohr, de Broglie, Einstein, Heisenberg, Planck, Rutherford, Schrödinger).
\end{itemize}

\subsection{Language and Literature Curriculum \mdseries(40 credits)}
\begin{itemize}[noitemsep, topsep=0pt]
    \item \textbf{Ancient Greek:} Four years of Ancient Greek: grammar, syntax, translation of Greek poetry and prose including Homer, Plato (Meno and other dialogues), Aristotle, Sophocles.
    \item \textbf{French:} Four years of French: grammar, style, translation (La Rochefoucauld, Molière, Racine, Baudelaire, modern poets).
    \item \textbf{English and Logic:} English grammar, style, composition. Formal logic based on Aristotle's Categories. Study of Chaucer and Shakespeare.
\end{itemize}

\subsection{Seminar: Philosophy, Literature, and Political Theory \mdseries(40 credits)}
Four-year discussion-based study of seminal texts in Western thought including works in philosophy, literature, political theory, economics, biblical literature, theology, psychology, and history. Approximately 50--150 pages per class meeting. Intensive preceptorial studies in junior and senior years, including Heidegger's \textit{Being and Time} (junior year) and \textit{Human Origins} (senior year). Annual essays and oral examinations each semester.

\subsection{Music \mdseries(12 credits)}
Elements of musical notation and sight reading. Melodic analysis, chorale counterpoint, and harmony. Close study of major works including choral music from Chant through four-part classics.

\end{document}
